A clean write-up is:

```tex
\begin{proof}
Assume for contradiction that
\[
|S\cap(S+v)|\ge 2.
\]
Choose distinct points
\[
p_1,p_2\in S\cap(S+v).
\]
Then
\[
p_1,p_2\in S,
\]
and because \(p_1,p_2\in S+v\), we also have
\[
p_1-v\in S
\qquad\text{and}\qquad
p_2-v\in S.
\]
Hence the four points
\[
p_1,\quad p_2,\quad p_1-v,\quad p_2-v
\]
all lie in \(S\).

Since \(|S|=3\), at least two of these four points must coincide. Now
\[
p_1\neq p_2
\]
by choice, and therefore
\[
p_1-v\neq p_2-v
\]
because translation preserves distinctness. Also,
\[
p_1\neq p_1-v
\qquad\text{and}\qquad
p_2\neq p_2-v
\]
since \(v\neq 0\).

So the only possible coincidences are the two cross-equalities
\[
p_1-v=p_2
\qquad\text{or}\qquad
p_2-v=p_1.
\]
By symmetry, it is enough to treat the case
\[
p_1-v=p_2.
\]
Then
\[
p_1=p_2+v.
\]
Since \(p_2\in S\cap(S+v)\), we also know that
\[
p_2-v\in S.
\]
Thus \(S\) contains the three points
\[
p_2-v,\qquad p_2,\qquad p_2+v=p_1.
\]

These three points are distinct. Indeed,
\[
p_2-v\neq p_2,\qquad p_2\neq p_2+v
\]
because \(v\neq 0\), and if \(p_2-v=p_2+v\), then \(2v=0\), hence \(v=0\), again impossible.

Moreover, all three points lie on the affine line
\[
p_2+\mathbb{R}v.
\]
So \(S\) consists of three distinct collinear points, contradicting that \(S=\{A,B,C\}\) is the vertex set of a non-degenerate triangle.

Therefore,
\[
|S\cap(S+v)|\le 1.
\]
\end{proof}

\begin{proof}[Conclusion for \(F\)]
Let
\[
F(p)=\mathbf{1}_S(p)-\mathbf{1}_{S+v}(p).
\]
Then \(F(p)\neq 0\) exactly when \(p\) belongs to exactly one of the sets \(S\) and \(S+v\). Hence
\[
\{p:F(p)\neq 0\}
=
\bigl(S\setminus(S+v)\bigr)\,\dot\cup\,\bigl((S+v)\setminus S\bigr),
\]
where the union is disjoint. Therefore
\[
|\{p:F(p)\neq 0\}|
=
|S\setminus(S+v)|+|(S+v)\setminus S|.
\]
Now
\[
|S\setminus(S+v)|=|S|-|S\cap(S+v)|
\]
and
\[
|(S+v)\setminus S|=|S+v|-|S\cap(S+v)|.
\]
Since \(|S|=|S+v|=3\), we get
\[
|\{p:F(p)\neq 0\}|
=
3+3-2|S\cap(S+v)|
=
6-2|S\cap(S+v)|.
\]
By the first proof,
\[
|S\cap(S+v)|\le 1,
\]
so
\[
|\{p:F(p)\neq 0\}|
\ge 6-2=4>3=|S|.
\]
Thus it is impossible in the two-translate case to have
\[
|\{p:F(p)\neq 0\}|=3=|S|.
\]
\end{proof}
```

The key contradiction is that two overlap points force the three vertices of \(S\) to be of the form
\[
q-v,\quad q,\quad q+v,
\]
which are collinear.